\section{Метод решения}

Прямой ход метода Гаусса за $n$ шагов приводит матрицу к треугольному виду.
На шаге $k$ выбирается ведущая строка с наибольшим $|a_{ik}|$ среди строк
$k, \ldots, n-1$ (выбор по столбцу уменьшает ошибки округления) и ставится на
место $k$. Затем из каждой строки ниже вычитается ведущая, умноженная на
$a_{ik}/a_{kk}$. Обратный ход находит $x_{n-1}, \ldots, x_0$. Прямой ход требует
около $\frac{2}{3}n^3$ операций, обратный --- $n^2$, поэтому распараллеливается
прямой ход.

Строки внутри шага независимы: каждая изменяется только своим потоком и читает
лишь ведущую строку. Строка $i$ обрабатывается потоком $i \bmod T$
(циклическое распределение): строки в работе убывают сверху, и при таком
распределении у всех потоков до конца остаётся почти одинаковое число строк.
Перестановка строк --- обмен двух указателей, данные не копируются.

Шаг состоит из двух частей, разделённых барьерами. Поток 0 выбирает и
переставляет ведущую строку; барьер; все потоки исключают свои строки; барьер.
Без первого барьера потоки читали бы ещё не выбранную ведущую строку, без
второго --- поток 0 переставлял бы строки, которые другие ещё изменяют.

Результат не зависит от числа потоков: каждый элемент вычисляется одними и теми
же операциями в том же порядке. Матрица считается вырожденной, если модуль
ведущего элемента не больше $n \cdot \varepsilon \cdot \max|a_{ij}|$, где
$\varepsilon$ --- машинная точность \texttt{double}.

\section{Описание программы}

\subsection{Модули}
\begin{itemize}
\item \texttt{common/}: коды возврата, вывод ошибок, разбор числовых ключей.
\item \texttt{os/}: команда потоков, барьер, монотонные часы, PID/TID, проверка терминала.
\item \texttt{src/gauss/system.c}: хранение, чтение, генерация системы, невязка.
\item \texttt{src/gauss/gauss.c}: прямой и обратный ход.
\item \texttt{src/gauss/main.c}: ключи запуска, замер времени, вывод результата.
\end{itemize}

Расширенная матрица хранится одним блоком памяти, доступ к строкам --- через
массив указателей. Система читается со стандартного ввода; если это терминал
(\texttt{isatty}), программа запрашивает размер и каждую строку. Для замеров
ключ \texttt{-g n} строит случайную матрицу с $b_i = \sum_j a_{ij}$, точное
решение которой --- все единицы. Правильность проверяется невязкой
$\max_i |(Ax)_i - b_i|$ по копии исходной системы.

\subsection{Потоки}

\texttt{os\_team\_run(T, body, context)} запускает тело в $T$ потоках сразу:
индекс 0 выполняет вызывающий (главный) поток, для остальных вызывается
\texttt{pthread\_create}. Поэтому при ключе \texttt{-t T} в процессе ровно $T$
потоков. Новые потоки сначала ждут у <<стартовых ворот>> (mutex и условная
переменная), пока не будут созданы все. Если \texttt{pthread\_create} завершится
ошибкой, ворота отменяются: созданные потоки выходят, не входя в барьер, который
уже никогда бы не заполнился. Все созданные потоки присоединяются через
\texttt{pthread\_join}.

\subsection{Синхронизация}

Шаги разделяет барьер \texttt{pthread\_barrier\_t}. Поток, пришедший к барьеру,
засыпает в ядре (\texttt{futex}); последний пришедший будит остальных. Барьер
также упорядочивает память: всё, что поток записал до барьера, видно другим
после него. Поэтому общие данные (строки, флаг вырожденности) не защищаются
блокировками: каждое поле в каждый момент пишет один поток, а читают его
другие только после барьера.

\subsection{Коды возврата}
\begin{center}
\begin{tabular}{ll}
0 & штатное завершение \\
2 & неверные ключи запуска \\
3 & ошибка чтения или формата системы \\
4 & ошибка записи решения \\
5 & матрица вырождена: нет единственного решения \\
6 & ошибка создания потока или барьера \\
7 & недостаточно памяти \\
\end{tabular}
\end{center}
